% ==============================================================================
% main.tex — DI41015 Computer Aided Engineering I: CFD Coursework Report
% Clean arXiv / NeurIPS (1706.03762) Pure Black-and-White Academic Manuscript
% Strict Academic Prose: Zero Code Fonts, Zero File Names, IEEE References
% ==============================================================================
\documentclass[10pt,a4paper]{article}
\usepackage[singlecol]{x57_paper}
\graphicspath{{../DELIVERABLES/figures/}{../DELIVERABLES/fluent_renders/}{assets/}}

% --- Title & Author (Compact NeurIPS / arXiv 1706.03762 Format) ---
\papertitle{Aerodynamic and Structural Analysis of the NASA X-57\\Distributed Electric Propulsion Aircraft}
\papersubtitle{Computer Aided Engineering I: CFD Coursework Report}
\paperauthor{Yuang Zhou}
\paperauthoremail{2617376@dundee.ac.uk}

\begin{document}
\makepapertitle

% --- Abstract intentionally left blank until all verified results are complete ---
\begin{abstractbox}
\vspace{16mm}
\end{abstractbox}

% ==============================================================================
\section{Introduction}
\label{sec:intro}
% ==============================================================================
Distributed electric propulsion (DEP) represents a major architectural concept for low-emission aviation, exemplified by the NASA X-57 Maxwell flight demonstrator~\cite{borer2016sceptor}. By distributing twelve compact electrically driven high-lift propellers (HLPs) along the leading edge of a high-aspect-ratio wing~\cite{patterson2016hlp} alongside two counter-rotating wingtip cruise propellers, DEP elevates local dynamic pressure across the lifting surface at low airspeeds~\cite{deere2017x57wing}—enabling a substantially smaller, lower-drag cruise wing while attenuating wingtip-vortex induced drag~\cite{deere2019modiii}.

During the onset of ground taxiing at sea level ($V_\infty \approx 0$), however, external aerodynamic behaviour differs markedly from forward cruise flight. In the absence of ambient freestream flow, the velocity and pressure fields are governed entirely by propeller momentum injection, near-field slipstream contraction, azimuthal swirl, and runway ground interaction. Evaluating take-off thrust, static wing lift increment, wing surface suction peaks, rotational speed scaling, high-lift design modifications, and structural safety margins therefore requires mesh-verified three-dimensional computational fluid dynamics (CFD) coupled with structural finite-element analysis.

The primary technical contributions of this investigation are summarised as follows:
\begin{enumerate}
  \item \textbf{Geometric Conditioning and Actuator-Disk Calibration}: The baseline three-dimensional aircraft geometry is restored to a level metric coordinate frame (\SI{0.006}{mm} spanwise symmetry tolerance and \SI{0.77853}{m} sponson clearance), and seven starboard inboard-up swirling actuator disks are calibrated against NASA reference propulsion maps with discrete momentum-integral discrepancies below $0.5\%$.
  \item \textbf{Systematic Poly-Hexcore Grid Verification}: A three-level systematic grid-refinement sequence (Coarse, Medium, and Fine) is constructed up to the solver cell limit ($< 1{,}048{,}576$ half-aircraft cells, corresponding to approximately $2.0 \times 10^6$ full-aircraft cells), with full six-step ASME V\&V~20 GCI uncertainty bounds and verified wall-resolved $y^+ \approx 1$ coverage over the blown wing.
  \item \textbf{Parametric Aeropropulsive Ablation Matrix}: Eighteen converged simulations systematically isolate propeller speed scaling, slipstream swirl, HLP-versus-wingtip aerodynamic coupling, alternating counter-rotation, runway ground effect, \SI{10}{m/s} take-off roll transition, and turbulence-closure sensitivity comparing SST $k$--$\omega$ against Realizable $k$--$\varepsilon$.
  \item \textbf{High-Lift Design Modifications}: Three-dimensional HLP $-10^\circ$ thrust vectoring and five high-resolution two-dimensional blown-aerofoil configurations (plain flap, slotted flap, leading-edge slat with slotted flap, and increased camber) are evaluated on a unified lift--drag Pareto frontier.
  \item \textbf{Fluid Shear and One-Way FSI Structural Convergence}: Computed surface pressure and wall-shear loads are mapped onto four-node structural shell finite-element models (SHELL181) across three structural mesh resolutions to verify converged von Mises stresses and global safety factors.
\end{enumerate}

\begin{figure}[t]
  \centering
  \IfFileExists{../DELIVERABLES/figures/fig22_3d_2d_hybrid_aeropropulsive_synthesis.pdf}{%
    \includegraphics[width=\textwidth]{fig22_3d_2d_hybrid_aeropropulsive_synthesis.pdf}%
    \caption{Overview of the aeropropulsive flow field and structural response of the NASA X-57 DEP aircraft under sea-level static taxiing: combined three-dimensional slipstream pathlines, $Q$-criterion vortex topology, converged one-way FSI structural stress distribution, and aerodynamic performance curves.}%
  }{%
    \includegraphics[width=\textwidth]{fig01_geometry_and_domain.pdf}%
    \caption{Restored level metric geometry of the NASA X-57 DEP aircraft and computational domain: (a)~plan view, (b)~front view showing the six starboard high-lift actuator disks (\SI{5400}{RPM}) and single wingtip cruise disk (\SI{2700}{RPM}), (c)~starboard elevation with measured \SI{0.77853}{m} sponson clearance, and (d)~external fluid domain extent and boundary conditions.}%
  }
  \label{fig:teaser}
\end{figure}

% ==============================================================================
\section{Geometry, Physical Modelling, and Numerical Methodology}
\label{sec:method}
% ==============================================================================
\subsection{Geometric Restoration, Reference Dimensions, and Domain Boundaries}
The supplied three-dimensional aircraft CAD geometry is provided in imperial units with an embedded attitude rotation relative to the level flight-attitude reference frame. Applying the inverse attitude transformation $\mathbf{R}^T$ and exact imperial-to-metric scaling ($\times \SI{0.0254}{m/in}$) establishes a level fuselage coordinate system ($+x$ aft, $+y$ vertical upward, $+z$ starboard) with a verified spanwise symmetry tolerance of \SI{0.006}{mm}. Removing the thirty-three static starboard propeller blades isolates five distinct aircraft aerodynamic surfaces—the main wing, fuselage, propulsion nacelles, main-landing-gear sponson, and empennage—characterised by a reference wing planform area $S_{\text{ref}} = \SI{6.19336}{m^2}$, full wingspan $b = \SI{9.63834}{m}$, and mean aerodynamic chord $\text{MAC} = \SI{0.64922}{m}$.

Exploiting spanwise symmetry about $z = 0$, the half-aircraft fluid domain extends over $x \in [-18, 30]\text{ m}$, $y \in [-1.70, 18]\text{ m}$ (with a stationary no-slip runway surface at $y = \SI{-1.70}{m}$), and $z \in [0, 20]\text{ m}$, as illustrated in Figure~\ref{fig:teaser}. All far-field external boundaries are prescribed as zero-gauge ambient pressure openings at sea-level standard atmospheric conditions ($P_{\text{op}} = \SI{101325}{Pa}$, $T_0 = \SI{288.15}{K}$).

\subsection{Governing Equations, Turbulence Closure, and Actuator-Disk Formulation}
Steady compressible ideal-gas RANS equations ($\mu = \SI{1.789e-5}{Pa.s}$, ambient reference density $\rho = \SI{1.225}{kg/m^3}$) are solved in Ansys Fluent 2025~R2~\cite{ansys2025fluent} employing the shear-stress transport (SST) $k$--$\omega$ turbulence model~\cite{menter1994sst} augmented with Spalart--Shur rotation and curvature correction~\cite{spalart1997sarc}, production limiting, and warped-face gradient correction. Because full blade-resolved sliding-mesh modelling of fourteen rotors is computationally prohibitive within the allowable mesh budget, each propeller is represented as a \SI{0.050}{m} thick swirling Goldstein-weighted actuator disk~\cite{adkins1994propellers}:
\begin{equation}
  S_x(r) = \frac{T}{\mathcal{I}_T}\,\hat{r}\sqrt{1-\hat{r}}, \qquad
  S_\theta(r) = \pm\frac{Q}{\mathcal{I}_Q}\,\frac{\sqrt{\hat{r}(1-\hat{r})}}{r}, \qquad
  \hat{r} = \frac{r - R_{\text{hub}}}{R_{\text{tip}} - R_{\text{hub}}},
  \label{eq:actuator}
\end{equation}
where $T$ and $Q$ denote the target rotor thrust and torque, $R_{\text{hub}}$ and $R_{\text{tip}}$ are the hub and tip radii, and the discrete cell-volume normalisation constants $\mathcal{I}_T$ and $\mathcal{I}_Q$ are evaluated directly on each production mesh to enforce discrete thrust and torque integration accuracy within $0.5\%$ of NASA X-57 blade-element momentum targets~\cite{nasa2024hlp,nasa2024cruise}.

\begin{pipelinebox}[Numerical Verification and Coupling Workflow]
\begin{enumerate}[leftmargin=1.5em, itemsep=1.5pt, topsep=2pt]
  \item \textbf{Geometric Conditioning and Rotor Calibration}: Restore the metric aircraft CAD surfaces, verify sub-millimetre sponson and trailing-edge clearances, and calibrate the seven inboard-up swirling actuator disks.
  \item \textbf{Systematic Three-Grid Poly-Hexcore Meshing}: Generate Coarse, Medium, and Fine conformal boundary-layer meshes targeting $y^+ \approx 1$ on the main wing surface, verifying positive cell volumes and solver minimum orthogonal quality exceeding $0.10$.
  \item \textbf{Steady RANS Verification and Parametric Matrix}: Solve the three-level grid-convergence sequence, rotational-speed sweep, swirl isolation, propulsion decomposition, ground-effect, take-off roll, turbulence sensitivity, and high-lift modification cases; compute six-step ASME V\&V~20 GCI bounds.
  \item \textbf{One-Way Fluid--Structure Structural Analysis}: Transfer converged aerodynamic wall pressure and wall shear stress fields onto four-node aluminium-alloy (Al-7075-T6) structural shell meshes across three refinement levels to verify von Mises stress and tip-deflection convergence.
\end{enumerate}
\end{pipelinebox}

% ==============================================================================
\section{Meshing Strategy, Quality Verification, and Grid Convergence}
\label{sec:mesh}
% ==============================================================================
A three-level systematic Poly-Hexcore mesh ladder (Coarse, Medium, and Fine) is constructed with conformal prism inflation layers on all aircraft walls and nested volumetric bodies of influence around the seven actuator disks and downstream slipstream wake. Detailed cell-count allocations, independent solver quality metrics, area-weighted wing $y^+$ statistics, and six-step ASME V\&V~20 GCI results~\cite{celik2008gci,asme2009vv20} will be populated in Table~\ref{tab:mesh_gci} upon completion of the grid-convergence simulations.

\begin{table}[t]
  \caption{Three-level Poly-Hexcore mesh audit and ASME V\&V~20 grid convergence summary at nominal 100\% rotational speed (swirl enabled).}
  \label{tab:mesh_gci}
  \centering
  \small
  \setlength{\tabcolsep}{5.5pt}
  \begin{tabular}{lccccccc}
    \toprule
    Grid Level & Half Cells & Full Equiv. & $\text{OQ}_{\min}$ & Wing $\bar{y}^+$ & $L_{\text{air}}$ (N) & $D_{\text{air}}$ (N) & $\text{GCI}_{\text{fine}}$ (\%) \\
    \midrule
    Coarse & --- & --- & --- & --- & --- & --- & --- \\
    Medium & --- & --- & --- & --- & --- & --- & --- \\
    \best{Fine} & --- & --- & --- & --- & --- & --- & --- \\
    \bottomrule
  \end{tabular}
\end{table}

% ==============================================================================
\section{Baseline 100\% RPM Aerodynamics, Pressure Field, and Pathlines}
\label{sec:baseline}
% ==============================================================================
This section will present the verified fine-grid baseline flow solution at nominal rotational speed (\SI{5400}{RPM} high-lift propellers and \SI{2700}{RPM} wingtip cruise propellers), covering gross and net static take-off thrust, component-wise aerodynamic force decomposition, wing chordwise pressure-coefficient distributions, exact spatial coordinates of the global minimum surface pressure, cross-sectional velocity and vorticity slices, and three-dimensional slipstream pathlines and $Q$-criterion vortex structures.

\begin{keyfinding}[Summary of Baseline Blown-Wing Suction Mechanism]
Verified quantitative findings regarding gross and net static thrust, upper- versus lower-wing pressure load balance, and the global peak-suction coordinates will be reported here upon completion of the fine-grid baseline solution.
\end{keyfinding}

% ==============================================================================
\section{Propeller Speed Scaling and Slipstream Swirl Physics}
\label{sec:rpm_swirl}
% ==============================================================================
This section will analyse the matched medium-grid rotational-speed sweep across 0\%, 50\%, 75\%, and 100\% nominal propeller speed to quantify the quadratic scaling collapse of aerodynamic forces and moments, followed by a direct comparison between swirl-enabled and swirl-disabled actuator-disk simulations to isolate axial vorticity transport and spanwise lift redistribution $\Delta(\mathrm{d}L/\mathrm{d}z)$.

% ==============================================================================
\section{Expanded Aeropropulsive Ablation, Ground Effect, and Model Sensitivity}
\label{sec:expanded}
% ==============================================================================
This section will present six additional medium-grid sensitivity and decomposition studies: independent high-lift-only versus wingtip-only propulsion decomposition to quantify nonlinear aeropropulsive synergy, alternating counter-rotating propeller directions, in-ground-effect versus out-of-ground-effect comparisons, \SI{10}{m/s} take-off ground-roll transition, and turbulence-closure sensitivity comparing SST $k$--$\omega$ against Realizable $k$--$\varepsilon$~\cite{shih1995rke}.

% ==============================================================================
\section{CFD-Validated High-Lift Design Modifications}
\label{sec:optimization}
% ==============================================================================
This section will evaluate aerodynamic modifications designed to increase static and low-speed lift, combining three-dimensional $-10^\circ$ high-lift thrust vectoring with five high-resolution two-dimensional blown-aerofoil configurations (clean baseline, $+15^\circ$ plain flap, $+30^\circ$ slotted flap, leading-edge slat with $+30^\circ$ slotted flap, and $+3\%$ camber increase) summarised on a unified $\Delta C_L$--$\Delta C_D$ Pareto frontier.

% ==============================================================================
\section{Fluid Shear Stresses and One-Way FSI Structural Convergence}
\label{sec:fsi}
% ==============================================================================
This section will examine both hydrodynamic stresses—surface wall shear stress $\tau_w$, modelled Reynolds shear stress $-\rho\overline{u'v'}$, and slipstream total-pressure loss—and the coupled one-way FSI structural response evaluated on three progressively refined four-node structural shell meshes (SHELL181~\cite{ansys2025shell181}), reporting von Mises equivalent stress convergence, global structural safety factors, root and pylon stress concentrations, and wing tip deflection.

% ==============================================================================
\section{Limitations, Reproducibility, and Conclusions}
\label{sec:conclusions}
% ==============================================================================
\subsection{Modelling Assumptions and Limitations}
Key physical approximations—including steady-state RANS time-averaging, actuator-disk omission of discrete blade-passing unsteadiness, solver cell-count constraints, and rigid-aerodynamic one-way structural coupling—will be critically assessed alongside the completed numerical results.

\subsection{Reproducibility and Visualisation Standards}
All numerical tables and figures employ perceptually uniform scientific colour maps~\cite{crameri2020colour} and consistent Times New Roman typography, derived directly from verified Ansys Fluent 2025~R2 and Ansys Mechanical APDL 2025~R2 simulations.

\subsection{Conclusions}
Concise quantitative conclusions addressing each aerodynamic, design-modification, and structural-verification objective will be finalised upon completion of the simulation matrix.

% ==============================================================================
\bibliographystyle{ieeetr}
\bibliography{references}

% ==============================================================================
\appendix
\section{Complete Parametric Matrix and Convergence Verification Summary}
Comprehensive aerodynamic performance metrics, force decompositions, and four-fold convergence diagnostics across all eighteen simulated cases will be tabulated in this appendix.

\section{ASME V\&V 20 GCI Formulation and Blade-Element Calibration}
The six-step ASME V\&V~20 discretisation uncertainty equations and per-rotor actuator-disk momentum-integral verification statistics will be documented in this appendix.

\section{Supplementary Three-Dimensional Flow and Structural Stress Visualisations}
Additional high-resolution three-dimensional visualisations of the fluid shear-stress fields, slipstream topology, and coupled structural stress distributions will be compiled in this appendix.

\end{document}
